% ========== 页面设置 ==========
\usepackage[top=2.5cm, bottom=2.5cm, left=3cm, right=3cm, headheight=24.5pt]{geometry}
\raggedbottom

% ========== 数学宏包 ==========
\usepackage{amsmath, amssymb, amsthm}
\usepackage{mathtools}
\usepackage{bm}           % 粗体数学符号
\usepackage{physics}      % 物理学常用符号

% ========== 图形宏包 ==========
\usepackage{graphicx}
\usepackage{caption}
\usepackage{tikz}
\usetikzlibrary{arrows.meta, calc, positioning, decorations.markings, decorations.pathmorphing, patterns, shapes.geometric}
\usepackage{pgfplots}
\pgfplotsset{compat=1.18}

% ========== 表格宏包 ==========
\usepackage{booktabs}
\usepackage{array}
\usepackage{multirow}

% ========== 其他宏包 ==========
\usepackage{hyperref}
\hypersetup{
    colorlinks=true,
    linkcolor=blue!70!black,      % 内部链接（目录、引用）
    urlcolor=blue!80!black,       % 外部URL
    citecolor=green!50!black      % 文献引用
}
\usepackage{xcolor}
\usepackage{tcolorbox}
\tcbuselibrary{skins, breakable, theorems}
\usepackage{enumitem}
\usepackage{fancyhdr}
\usepackage{emptypage}
\usepackage{titlesec}
\usepackage{indentfirst}
\usepackage{float}
\usepackage{needspace}
\usepackage{pifont}           % 特殊符号(如\ding)
\usepackage{makeidx}          % 索引
\makeindex

% ========== 颜色系统 ==========
% 设计原则：主色调（蓝）+ 辅助色（青绿、紫、棕）+ 警示色（红）

% --- 主色调：深蓝 ---
\definecolor{theoremblue}{RGB}{25, 84, 140}       % 定理、标题（沉稳深蓝）
\definecolor{mainblue}{RGB}{25, 84, 140}           % 别名，用于标题

% --- 辅助色 ---
\definecolor{defteal}{RGB}{0, 128, 128}            % 定义（青绿色，区别于蓝但协调）
\definecolor{examplepurple}{RGB}{106, 90, 140}     % 例题（柔和薰衣紫）
\definecolor{notebrown}{RGB}{160, 100, 60}         % 注、思考题（温暖琥珀棕）
\definecolor{hintteal}{RGB}{70, 150, 150}          % 提示框（清新青色）
\definecolor{physicsbrown}{RGB}{139, 90, 43}       % 物理背景框（深琥珀）

% --- 警示色 ---
\definecolor{alertred}{RGB}{180, 50, 50}          % 警告、易错点

% --- 公式与高亮 ---
\definecolor{formulabg}{RGB}{255, 252, 245}        % 公式框背景（暖白）
\definecolor{formulaframe}{RGB}{180, 140, 80}      % 公式框边框（金色）
\definecolor{keyblue}{RGB}{30, 100, 180}           % 关键词蓝色
\definecolor{importantred}{RGB}{200, 50, 50}       % 重要内容红色
\colorlet{chapteraccent}{formulaframe!88!black}

% --- 封面与全书可复用的补充色 ---
\definecolor{coverpaper}{RGB}{248, 246, 240}
\definecolor{covernavy}{RGB}{26, 48, 88}
\definecolor{covergold}{RGB}{186, 146, 76}
\definecolor{coverline}{RGB}{196, 202, 212}
\definecolor{coverbrick}{RGB}{174, 96, 74}
\definecolor{coverteal}{RGB}{82, 132, 136}
\definecolor{coverplum}{RGB}{104, 97, 142}

% ========== 配图统一样式 ==========
\colorlet{figaxis}{black!78}
\colorlet{figguide}{black!28}
\colorlet{figblue}{mainblue!92!black}
\colorlet{figteal}{defteal!88!black}
\colorlet{figpurple}{examplepurple!92!black}
\colorlet{figred}{alertred!88!black}
\colorlet{figbrown}{physicsbrown!88!black}
\colorlet{figgold}{notebrown!85!black}
\colorlet{figsoftfill}{black!2}
\colorlet{figbluefill}{figblue!7}
\colorlet{figtealfill}{figteal!7}
\colorlet{figpurplefill}{figpurple!7}
\colorlet{figredfill}{figred!7}
\colorlet{figbrownfill}{figbrown!12}

\tikzset{
  mp axis/.style={-{Stealth[length=2.6mm,width=1.8mm]}, line width=0.85pt, draw=figaxis},
  mp thin axis/.style={-{Stealth[length=2.2mm,width=1.5mm]}, line width=0.75pt, draw=figaxis},
  mp helper/.style={draw=figguide, dash pattern=on 2.2pt off 1.8pt, line width=0.55pt},
  mp soft/.style={draw=black!18, line width=0.5pt},
  mp circle/.style={draw=figblue, line width=0.9pt},
  mp angle/.style={draw=figbrown, line width=0.9pt},
  mp vector/.style={-{Stealth[length=2.7mm,width=1.8mm]}, line width=1.05pt, draw=figblue},
  mp vector b/.style={-{Stealth[length=2.7mm,width=1.8mm]}, line width=1.05pt, draw=figteal},
  mp vector c/.style={-{Stealth[length=2.7mm,width=1.8mm]}, line width=1.1pt, draw=figpurple},
  mp vector d/.style={-{Stealth[length=2.7mm,width=1.8mm]}, line width=1.05pt, draw=figred},
  mp vector b dashed/.style={mp vector b, dash pattern=on 3.4pt off 2.1pt},
  mp vector d dashed/.style={mp vector d, dash pattern=on 3.2pt off 2.0pt},
  mp point/.style={circle, fill=figblue, inner sep=1.7pt},
  mp point b/.style={circle, fill=figteal, inner sep=1.7pt},
  mp point c/.style={circle, fill=figpurple, inner sep=1.7pt},
  mp point d/.style={circle, fill=figred, inner sep=1.7pt},
  mp area blue/.style={fill=figbluefill, draw=none},
  mp area teal/.style={fill=figtealfill, draw=none},
  mp area purple/.style={fill=figpurplefill, draw=none},
  mp area red/.style={fill=figredfill, draw=none},
  mp label/.style={font=\small},
  mp small/.style={font=\footnotesize},
  mp card/.style={
    rounded corners=2pt,
    draw=black!16,
    fill=figsoftfill,
    line width=0.5pt,
    inner sep=4pt,
    align=center
  }
}

% ========== 图题与列表 ==========
\captionsetup{
  font=small,
  labelfont={bf, color=theoremblue},
  textfont={small},
  labelsep=space,
  skip=6pt
}

\setlist[itemize]{
  leftmargin=2.1em,
  itemsep=0.25em,
  topsep=0.4em,
  parsep=0pt
}

\setlist[enumerate]{
  leftmargin=2.4em,
  itemsep=0.3em,
  topsep=0.45em,
  parsep=0pt
}

% ========== 定理环境（使用 newtcbtheorem，标题独立色块）==========

% 定理环境（蓝色）
\newtcbtheorem[number within=section]{dingli}{定理}{
  enhanced,
  colback=theoremblue!5!white,
  colframe=theoremblue,
  coltitle=white,
  fonttitle=\bfseries,
  separator sign none,
  boxrule=1pt,
  arc=2pt,
  left=8pt, right=8pt, top=6pt, bottom=6pt,
  attach boxed title to top left={yshift=-2mm, xshift=5mm},
  boxed title style={
    colback=theoremblue,
    colframe=theoremblue,
    boxrule=0pt,
    arc=2pt
  },
  breakable,
  before skip=0.8em,
  after skip=0.8em
}{dingli}

% 定义环境（青绿色）
\newtcbtheorem[number within=section]{dingliyi}{定义}{
  enhanced,
  colback=defteal!5!white,
  colframe=defteal,
  coltitle=white,
  fonttitle=\bfseries,
  separator sign none,
  boxrule=1pt,
  arc=2pt,
  left=8pt, right=8pt, top=6pt, bottom=6pt,
  attach boxed title to top left={yshift=-2mm, xshift=5mm},
  boxed title style={
    colback=defteal,
    colframe=defteal,
    boxrule=0pt,
    arc=2pt
  },
  breakable,
  before skip=0.8em,
  after skip=0.8em
}{dingliyi}

% 引理环境（蓝色）
\newtcbtheorem[number within=section]{yinli}{引理}{
  enhanced,
  colback=theoremblue!5!white,
  colframe=theoremblue,
  coltitle=white,
  fonttitle=\bfseries,
  separator sign none,
  boxrule=1pt,
  arc=2pt,
  left=8pt, right=8pt, top=6pt, bottom=6pt,
  attach boxed title to top left={yshift=-2mm, xshift=5mm},
  boxed title style={
    colback=theoremblue,
    colframe=theoremblue,
    boxrule=0pt,
    arc=2pt
  },
  breakable,
  before skip=0.8em,
  after skip=0.8em
}{yinli}

% 推论环境（蓝色）
\newtcbtheorem[number within=section]{tuilun}{推论}{
  enhanced,
  colback=theoremblue!5!white,
  colframe=theoremblue,
  coltitle=white,
  fonttitle=\bfseries,
  separator sign none,
  boxrule=1pt,
  arc=2pt,
  left=8pt, right=8pt, top=6pt, bottom=6pt,
  attach boxed title to top left={yshift=-2mm, xshift=5mm},
  boxed title style={
    colback=theoremblue,
    colframe=theoremblue,
    boxrule=0pt,
    arc=2pt
  },
  breakable,
  before skip=0.8em,
  after skip=0.8em
}{tuilun}

% 例题环境（紫色）
\newtcbtheorem[number within=section]{lizi}{例题}{
  enhanced,
  colback=examplepurple!5!white,
  colframe=examplepurple,
  coltitle=white,
  fonttitle=\bfseries,
  separator sign none,
  boxrule=1pt,
  arc=2pt,
  left=8pt, right=8pt, top=6pt, bottom=6pt,
  attach boxed title to top left={yshift=-2mm, xshift=5mm},
  boxed title style={
    colback=examplepurple,
    colframe=examplepurple,
    boxrule=0pt,
    arc=2pt
  },
  breakable,
  before skip=0.8em,
  after skip=0.8em
}{lizi}

% 注环境（琥珀棕色）
\newtcbtheorem[number within=section]{zhu}{注}{
  enhanced,
  colback=notebrown!5!white,
  colframe=notebrown,
  coltitle=white,
  fonttitle=\bfseries,
  separator sign none,
  boxrule=1pt,
  arc=2pt,
  left=8pt, right=8pt, top=6pt, bottom=6pt,
  attach boxed title to top left={yshift=-2mm, xshift=5mm},
  boxed title style={
    colback=notebrown,
    colframe=notebrown,
    boxrule=0pt,
    arc=2pt
  },
  breakable,
  before skip=0.6em,
  after skip=0.6em
}{zhu}

% 思考题环境（琥珀棕色）
\newtcbtheorem[number within=section]{sikao}{思考题}{
  enhanced,
  colback=notebrown!5!white,
  colframe=notebrown,
  coltitle=white,
  fonttitle=\bfseries,
  separator sign none,
  boxrule=1pt,
  arc=2pt,
  left=8pt, right=8pt, top=6pt, bottom=6pt,
  attach boxed title to top left={yshift=-2mm, xshift=5mm},
  boxed title style={
    colback=notebrown,
    colframe=notebrown,
    boxrule=0pt,
    arc=2pt
  },
  breakable,
  before skip=0.6em,
  after skip=0.6em
}{sikao}



% 证明环境（灰色边框，区别于定理）
% 重新定义 proof 环境，带有浅色背景
\renewenvironment{proof}{
  \begin{tcolorbox}[
    enhanced,
    colback=gray!5!white,
    colframe=gray!60!black,
    boxrule=0pt,
    leftrule=2pt,
    arc=0pt,
    left=8pt, right=6pt, top=4pt, bottom=4pt,
    breakable,
    before skip=0.5em,
    after skip=0.5em
  ]
  \textbf{证明：}\par\vspace{2pt}
}{
  \par\vspace{2pt}\hfill$\blacksquare$
  \end{tcolorbox}
}



% ========== 彩色公式框 ==========
% 重要公式框（金色边框）
\newtcolorbox{zhongyao}[1][]{
  enhanced,
  colback=formulabg,
  colframe=formulaframe,
  coltitle=white,
  fonttitle=\bfseries,
  boxrule=1.5pt,
  arc=3pt,
  title=#1,
  attach boxed title to top center={yshift=-2mm},
  boxed title style={
    colback=formulaframe!85!black,
    colframe=formulaframe!85!black,
    boxrule=0pt,
    arc=2pt
  },
  before upper={\raggedright},
  after upper={\par},
  breakable
}

% 关键词高亮命令
\newcommand{\keyword}[1]{\textcolor{keyblue}{\textbf{#1}}}
\newcommand{\important}[1]{\textcolor{importantred}{\textbf{#1}}}
\newcommand{\formulaf}[1]{\boxed{\textcolor{theoremblue}{#1}}}

% 解答环境（仿 proof 的行文式环境，避免盒子套盒子带来的分页留白）
\makeatletter
\newenvironment{jie}[1][解]{%
  \par
  \pushQED{\qed}%
  \normalfont
  \topsep=0.35em\partopsep=0pt\relax
  \trivlist
  \item[\hskip\labelsep\textcolor{examplepurple}{\bfseries #1：}]%
  \ignorespaces
}{%
  \popQED\endtrivlist\@endpefalse
}
\makeatother

% ========== 自定义环境 ==========
% 警示框：常见错误（警示红）
\newtcolorbox{jinshi}[1][]{
  enhanced,
  colback=alertred!5!white,
  colframe=alertred,
  coltitle=white,
  fonttitle=\bfseries,
  title=#1,
  boxrule=0.9pt,
  arc=3pt,
  left=8pt, right=8pt, top=6pt, bottom=6pt,
  attach boxed title to top left={yshift=-2mm, xshift=5mm},
  boxed title style={
    colback=alertred!85!black,
    colframe=alertred!85!black,
    boxrule=0pt,
    arc=2pt
  },
  breakable,
  before skip=0.7em,
  after skip=0.8em
}

% 提示框（清新青色）
\newtcolorbox{tishi}[1][]{
  enhanced,
  colback=hintteal!5!white,
  colframe=hintteal,
  coltitle=white,
  fonttitle=\bfseries,
  title=#1,
  boxrule=0.9pt,
  arc=3pt,
  left=8pt, right=8pt, top=6pt, bottom=6pt,
  attach boxed title to top left={yshift=-2mm, xshift=5mm},
  boxed title style={
    colback=hintteal!85!black,
    colframe=hintteal!85!black,
    boxrule=0pt,
    arc=2pt
  },
  before upper={\setlength{\parindent}{2em}},
  breakable,
  before skip=0.75em,
  after skip=0.8em
}

% 物理背景框（深琥珀）
\newtcolorbox{wulibj}[1][]{
  enhanced,
  colback=physicsbrown!5!white,
  colframe=physicsbrown,
  coltitle=white,
  fonttitle=\bfseries,
  title=#1,
  boxrule=0.9pt,
  arc=3pt,
  left=8pt, right=8pt, top=6pt, bottom=6pt,
  attach boxed title to top left={yshift=-2mm, xshift=5mm},
  boxed title style={
    colback=physicsbrown!85!black,
    colframe=physicsbrown!85!black,
    boxrule=0pt,
    arc=2pt
  },
  breakable,
  before skip=0.8em,
  after skip=0.85em
}

% 名言框（用于历史引言和名人语录）
% 简洁引号式：无背景框 + 引号包裹 + 楷体深蓝灰
\newtcolorbox{mingyan}[1][]{
  enhanced,
  colback=white,
  colframe=white,
  boxrule=0pt,
  arc=0pt,
  left=0pt, right=0pt, top=4pt, bottom=4pt,
  before upper={``\kaishu\color{gray!60!blue}},
  after upper={''},
  breakable
}

% ========== 标题格式 ==========
\titleformat{\part}[display]
  {\thispagestyle{empty}\centering\Huge\bfseries\color{theoremblue}}
  {第\chinese{part}部分}
  {0.9em}
  {}
\titlespacing*{\part}{0pt}{0pt}{0pt}

\titleformat{\chapter}[display]
  {\centering}
  {\fontsize{19pt}{23pt}\selectfont\bfseries\color{theoremblue!82!black}第\chinese{chapter}章}
  {0.75em}
  {\fontsize{25pt}{31pt}\selectfont\bfseries\color{theoremblue}}
  [
    \vspace{0.7em}
    {\color{chapteraccent}\rule{0.24\textwidth}{1.2pt}}
  ]
\titlespacing*{\chapter}{0pt}{-14pt}{22pt}



\titleformat{\section}[block]
  {\Large\bfseries\color{theoremblue!80!black}}
  {\thesection}
  {0.8em}
  {}
\titlespacing*{\section}{0pt}{2.3ex plus 0.7ex minus 0.2ex}{1.1ex}


\titleformat{\subsection}[block]
  {\large\bfseries\color{mainblue!80!black}}
  {\thesubsection}
  {0.6em}
  {}
\titlespacing*{\subsection}{0pt}{1.9ex plus 0.6ex minus 0.2ex}{0.8ex}


% ========== 页眉页脚 ==========
\pagestyle{fancy}
\fancyhf{}
\fancyhead[LE, RO]{\small\color{black!62} 数学物理方法}
\fancyhead[RE, LO]{\small\color{black!62}\nouppercase{\leftmark}}
\fancyfoot[CE, CO]{\color{theoremblue!82!black}\thepage}
\renewcommand{\headrulewidth}{0.45pt}
\renewcommand{\headrule}{\hbox to\headwidth{\color{black!18}\leaders\hrule height \headrulewidth\hfill}}
